%==============================================================================
% Research Methods - Group assignment - proposals
%==============================================================================

% This loads the article template. For a text written in English, add the
% "english" option, i.e. \documentclass[english]{hogent-article}
\documentclass{hogent-article}

% Include bibliography file
\addbibresource{references.bib}

% Info about the assignment, course and study programme

\studyprogramme{Professionele bachelor toegepaste informatica}
% TODO: replace 2N1 with the correct class group(s)!
\course{Research Methods, 2N1} 
\assignmenttype{Groepsopdracht}
\academicyear{2025-2026}
\title{Ontwerp en evaluatie van een digitale ecologische voetafdruk-tracker ter ondersteuning van duurzaam dagelijkse gedrag bij Vlaamse scholieren}

% TODO: English text? Comment or remove the lines above and uncomment the lines
% below

% \studyprogramme{Professional bachelor in applied informatics}
% \course{Research Methods, INT-IT/2B}
% \assignmenttype{Group assignment}
% \academicyear{2025-2026}
% \title{Written preparation of research proposals}

% TODO: Replace student names and email addresses
\author{Ewout Van Mullem}
\email{ewout.vanmullem@student.hogent.be}

\author{Manu  Coppens}
\email{manu.coppens@student.hogent.be}

\author{Peter  Cypers}
\email{peter.cypers@student.hogent.be}

\author{Stef Simons}
\email{stef.simons@student.hogent.be}

% TODO: Link your Github-repository here
\projectrepo{https://github.com/HoGentTIN/rm-25-26-groep-85}

% TODO: Choose the appropriate specialization from the list below, 
%
% - Mobile \& Enterprise development
% - AI \& Data Engineering
% - Functional \& Business Analysis
% - System \& Network Engineering
% - Mainframe Expert
%
% - If your research does not fit into any of these domains, specify a specialization yourself.
%
\specialisation{Mobile \& Enterprise development}

% TODO: choose some appropriate keywords for your research proposals
\keywords{Carabon footprint, Ecological footprint, Ecological footprint calculator, User experience, Sustainable consumption behavior, Sustainability}

\begin{document}

% TODO: Replace the dummy text below with the actual contents of your research
% proposal. **ALSO REPLACE THE HEADING TITLES BY APPROPRIATE ONES!** The
% current headings are just placeholders.
\begin{abstract}
    %todo
  Wij presenteren het éérste onderzoek naar radiofrequente interferentie (RFI) op de toekomstige locatie van de laagfrequente Square Kilometre Array (SKA), het Murchison Radio-astronomie Observatorium (MRO), dat de RFI zowel in de tijd als ruimtelijk oplost.  Het onderzoek wordt uitgevoerd in een frequentiebereik van 1 MHz binnen de FM-band, ontworpen om de dichtstbijzijnde en sterkste FM-zenders van de MRO (gevestigd in Geraldton, op ongeveer 300 km afstand) te omvatten.  Gedurende ongeveer drie dagen, met gebruikmaking van de tweede iteratie van de Engineering Development Array in een ``all-sky imaging mode'', hebben we een reeks RFI-signalen gevonden.  We kunnen de signalen categoriseren in: signalen die rechtstreeks van de zenders worden ontvangen, vanaf hun horizonlocaties; reflecties van vliegtuigen (die ongeveer 13\% van de waarnemingsduur in beslag nemen); reflecties van objecten in een baan om de aarde; en reflecties van meteoor-ionisatiesporen.  In totaal analyseren we 33.994 beelden met een tijdresolutie van 7,92 s in beide polarisaties en een hoekresolutie van ongeveer 3,5$^{\circ}$, waarbij ongeveer veertigduizend RFI-gebeurtenissen worden gedetecteerd.  Deze gedetailleerde uitsplitsing van RFI in de MRO-omgeving zal toekomstige gedetailleerde analyses mogelijk maken van de waarschijnlijke gevolgen van RFI voor belangrijke wetenschap bij lage radiofrequenties met de SKA.
\end{abstract}

\tableofcontents

\section{Inleiding}%
\label{sec:inleiding}
%todo
De laatste jaren is veel onderzoek gewijd aan de uitrol van het Internet; helaas hebben slechts weinigen onderzoek gedaan naar de simulatie van wide area netwerken~\autocite{Hykes2013}. In deze position paper ont\-krach\-ten we het begrip van het World Wide Web. De notie dat theoretici meewerken aan de verbetering van gerandomiseerde algoritmen wordt meestal belangrijk geacht. De analyse van lambda calculus zou de verfijning van het World Wide Web enorm vergroten.

\begin{align*}
  f(x) &= x^2\\
  \mathbf{ g(x) } &= \frac{1}{x}\\
   F(x)  &= \int^a_b \frac{1}{3}x^3\\
\end{align*}

Wij bevestigen dat het veel aangeprezen certificeerbare algoritme voor de constructie van online algoritmen door Lee en Davis loopt in $\Theta(n^2)$ tijd. Het lijkt op het eerste gezicht pervers maar het viel in met onze verwachtingen. Bestaande lossless en coöperatieve heuristieken gebruiken superblocks om DHCP in te zetten. Maar, twee eigenschappen maken deze oplossing perfect: YnowHip simuleert alomtegenwoordige symmetrieën, en bovendien biedt YnowHip gerepliceerde symmetrieën. Deze combinatie van eigenschappen is nog niet verbeterd in eerder werk.

\section{Literatuurstudie}%
\label{sec:literatuurstudie}
% bron: Kok2021
Inleidende deel $\ldots$
%todo

%    blz.3
%    "In summary, limited findings suggest that ecological footprint calculators may be able to play a role in changing behaviour, but at the very least, this would require adoption and (perhaps repeated) use of the tool, which is currently difficult to achieve. Furthermore, little is known about the (positive and negative, intended and unintended) effects calculators may have on their users, on a psychological, behavioural, and social level. We thus argue that the final two steps of the IM approach have not been successfully executed for ecological footprint calculators."


\subsection{Carbon Footprint}
\label{ssec:carbonfootprint}
%disambiguation: wat is dit? wat houdt het in? welke acties van mensen hebben hier een invloed op?
%todo

\subsection{Tools om de carbon footprint te meten}
\label{ssec:tools}
% dit ssec kan misschien weggelaten worden...
%todo


\subsection{Impact van menselijke gewoontes op het klimaat}
% welke menselijke gewoontes hebben een positief/negatief invloed op het klimaat?

\subsection{Design \& bestaande apps}
\label{ssec:designenbestaandeapps}
% bron: Petersen2020 (Smiling Earth)
% https://www.wwf.nl/kom-in-actie/voetafdruktest
%todo

\subsection{Gebruiker-retentie en gamification}
\label{ssec:retentieengamification}
% bron: Oliveira2024
% Fogg's Behavior Model
% octalysis framework
% PBL-framework

\section{Methodologie}%
\label{sec:methodologie}

\subsection{Platformkeuze}
\label{ssec:platformkeuze}
% Mobile of Web-app -> een onderbouwde keuze maken
%todo

\subsection{Dagelijkse gewoontes met een invloed op het klimaat}
% uitzoeken welke dagelijkse gewoontes een impact hebben op het klimaat en belonen van positief gedrag met een puntensysteem

\subsubsection{Meetbare waarden}
% zijn er bepaalde meetbare waarden waarmee we bepaalde grafieken kunnen opstellen, of andere visualisaties?
% mogelijks zijn er bepaalde gewoontes die geen exacte meeteenheid hebben of enkel een ja/neen uitkomst hebben (vandaag niet met de auto gereden)(iets gerecycleerd)

\subsection{Gamification (implementatie)}
% welke gamification technieken we willen toepassen in onze app (gebaseerd op het literatuurstudie)
% hoe gaan we ze implementeren?

%todo Hoe meten van success? 


\section{Verwacht resultaat en conclusie}%
\label{sec:results}
% technisch: proof of concept, wireframes/mockups of een werkend app/website
% onderzoekend: is nog niet zeker wat als resultaat van het onderzoekende luik zal kunnen geconcludeerd worden...
% wat is het gevonden probleem? hoe wordt het opgelost en welke conclusies trekken we hieruit?
%todo


\printbibliography[heading=bibintoc]

\end{document}